\documentclass[14pt, a4paper]{article}
\usepackage{minitoc}
\usepackage[left=3.00cm, right=2.5cm, top=2.00cm, bottom=2.00cm]{geometry}
\usepackage{amsmath}
\usepackage{amssymb}
\usepackage{amsthm}
\usepackage{mathtools}
\usepackage{graphicx}
%\usepackage{algpseudocode}
%\usepackage{algorithm}
\usepackage[ruled,vlined,linesnumbered]{algorithm2e}
\usepackage{blindtext}
\usepackage{setspace}
\usepackage[utf8]{inputenc}
\usepackage[utf8]{vietnam}
\usepackage[center]{caption}
\usepackage[shortlabels]{enumitem}
\usepackage{fancyhdr} % header, footer
\usepackage{hyperref} % loại bỏ border với mục lục và công thức
\usepackage[nonumberlist, nopostdot, nogroupskip]{glossaries}
\usepackage{glossary-superragged}
\usepackage{tikz,tkz-tab}
\usepackage{pythonhighlight}
\setglossarystyle{superraggedheaderborder}
\pagestyle{fancy}
%\usepackage[style=numeric,sortcites]{biblatex}
%\addbibresource{ref.bib}
%\usepackage[numbers]{natbib}
\usepackage{indentfirst}
\usepackage[natbib,backend=biber,style=ieee, sorting=ynt]{biblatex}
\bibliography{ref.bib}

\graphicspath{{./figures/}}

\fancyhf{}
\rhead{\textbf{Môn học: Toán rời rạc và thuật toán}}
\lhead{\textbf{GVHD: PGS. TS. Nguyễn Thị Hồng Minh}}
\rfoot{\thepage}
\lfoot{\textbf{Nhóm học viên thực hiện: Nhóm 01 - Cao học Khoa học dữ liệu - K4}}
\renewcommand{\headrulewidth}{0.4pt}
\renewcommand{\footrulewidth}{0.4pt}
%
%\numberwithin{equation}{section}
%\numberwithin{algorithm}{section}
%\numberwithin{figure}{section}
%
%\setlength{\parindent}{0.5cm}
%
%\setcounter{secnumdepth}{3} % Cho phép subsubsection trong report
%\setcounter{tocdepth}{3} % Chèn subsubsection vào bảng mục lục

%\newtheorem{dl}{Định lý}
%\newtheorem{md}{Mệnh đề}
%\newtheorem{bd}{Bổ đề}
%\newtheorem{dn}{Định nghĩa}
%\newtheorem{hq}{Hệ quả}

%\newtheorem{baitap}{Bài tập}
%\newtheorem*{loigiai}{Lời giải}

%\numberwithin{dl}{section}
%\numberwithin{md}{section}
%\numberwithin{bd}{section}
%\numberwithin{dn}{section}
%\numberwithin{hq}{section}

\setlength{\parindent}{0cm}

\newtheorem{dl}{Định lý}
\newtheorem{bd}{Bổ đề}
\newtheoremstyle{sltheorem}
{}                % Space above
{}                % Space below
{\normalfont}        % Theorem body font % (default is "\upshape")
{}                % Indent amount
{\bfseries}       % Theorem head font % (default is \mdseries)
{.}               % Punctuation after theorem head % default: no punctuation
{ }               % Space after theorem head
{}                % Theorem head spec
\theoremstyle{sltheorem}
\newtheorem{baitap}{Bài tập}
\newtheoremstyle{soltheorem}
{}                % Space above
{}                % Space below
{\normalfont}        % Theorem body font % (default is "\upshape")
{}                % Indent amount
{\bfseries}       % Theorem head font % (default is \mdseries)
{.}               % Punctuation after theorem head % default: no punctuation
{\newline}               % Space after theorem head
{}                % Theorem head spec
\theoremstyle{soltheorem}
\newtheorem*{loigiai}{Lời giải}

\onehalfspacing


\begin{document}

    \begin{titlepage}

        \newcommand{\HRule}{\rule{\linewidth}{0.5mm}} % Defines a new command for the horizontal lines, change thickness here

        \center % Center everything on the page

        %----------------------------------------------------------------------------------------
        %	HEADING SECTIONS
        %----------------------------------------------------------------------------------------
        \textsc{\LARGE Đại học Quốc Gia Hà Nội}\\[0.5cm]
        \textsc{\LARGE Trường đại học Khoa học tự nhiên}\\[0.5cm] % Name of your university/college
        \textsc{\LARGE Khoa Toán - Cơ - Tin học}\\[0.5cm]

        \includegraphics[scale=0.2]{HUS-logo.jpg}\\[0.5cm]

        \textsc{\Large Chuyên ngành: Khoa học dữ liệu}\\[0.5cm] % Major heading such as course name


        %----------------------------------------------------------------------------------------
        %	TITLE SECTION
        %----------------------------------------------------------------------------------------

        \HRule \\[0.4cm]
        { \huge \bfseries Bài tập môn học}\\[0.4cm] % Title of your document
        \HRule \\[1.5cm]

        \textsc{\Large Môn học: Toán rời rạc và thuật toán}\\[1cm] % Minor heading such as course title


        \textsc{\Large Bài tập 2: Đường đi trong đồ thị và ứng dụng}\\[1cm]


        %----------------------------------------------------------------------------------------
        %	AUTHOR SECTION
        %----------------------------------------------------------------------------------------
        \begin{minipage}{0.4\textwidth}
            \begin{flushleft} \large
            \emph{Giảng viên hướng dẫn:} \\
            PGS. TS. Nguyễn Thị Hồng Minh % Supervisor's Name
            \end{flushleft}
        \end{minipage}\\[0.5cm]

        \begin{minipage}{0.4\textwidth}
        \begin{flushleft} \large
        \emph{Nhóm học viên thực hiện:}\\
        Nguyễn Chí Thanh \\
        MSHV: 21007925 \\ % Your name
        Vũ Ngọc Đại \\
        MSHV: 21007977 \\
        Vũ Minh Hưng \\
        MSHV: 21007973 \\
        Lê Diệu Thúy \\
        MSHV: 21007922 \\
        Lớp: Khoa học dữ liệu - K4
        \end{flushleft}
        \end{minipage}


        % If you don't want a supervisor, uncomment the two lines below and remove the section above
        %\Large \emph{Author:}\\
        %John \textsc{Smith}\\[3cm] % Your name

        %----------------------------------------------------------------------------------------
        %	DATE SECTION
        %----------------------------------------------------------------------------------------

        % I don't want day because it is English
        % {\large \today}\\[2cm] % Date, change the \today to a set date if you want to be precise

        %----------------------------------------------------------------------------------------
        %	LOGO SECTION
        %----------------------------------------------------------------------------------------

        %\includegraphics{logo/rsz_3logo-khtn.png}\\[1cm] % Include a department/university logo - this will require the graphicx package

        %----------------------------------------------------------------------------------------

        \vfill % Fill the rest of the page with whitespace

    \end{titlepage}

    \nocite{*}

    \newpage

    \begin{baitap}
        Hãy trình bày nhận thức, hiểu biết về đồ thị (graph)
        \begin{itemize}
            \item Những khái niệm cơ bản
            \item Cấu trúc dữ liệu biểu diễn đồ thị
            \item Mô hình và ứng dụng của đồ thị trong bài toán thực tiễn
        \end{itemize}
    \end{baitap}

    \begin{loigiai}

        Hãy trình bày nhận thức, hiểu biết về đồ thị (graph)

        \begin{itemize}
            \item Những khái niệm cơ bản:
        \end{itemize}
    \end{loigiai}

    \begin{baitap}
        Bài toán tìm đường đi (ngắn nhất) trên đồ thị và ứng dụng

        \begin{itemize}
            \item Bài toán
            \item Một số phương pháp (thuật toán) giải bài toán tìm đường đi (ngắn nhất)
            \item Ứng dụng: ví dụ một bài toán có thể mô hình hóa bằng đồ thị và giải bằng
            thuật toán tìm đường đi (ngắn nhất).
            \begin{itemize}
                \item Phát biểu bài toán
                \item Thuật toán
                \item Cài đặt chương trình.
            \end{itemize}
        \end{itemize}
    \end{baitap}

    \begin{loigiai}

        Bài toán tìm đường đi (ngắn nhất) trên đồ thị và ứng dụng

        \begin{itemize}
            \item Bài toán tìm đường đi (ngắn nhất) trên đồ thị:
            
            Trong các ứng dụng thực tế, chẳng hạn trong mạng lưới giao thông đường bộ, đường thuỷ hoặc
            đường không. Người ta không chỉ quan tâm đến việc tìm đường đi giữa hai địa điểm mà còn phải
            lựa chọn một hành trình tiết kiệm nhất (theo tiêu chuẩn không gian, thời gian hay chi phí). Khi đó
            phát sinh yêu cầu tìm đường đi ngắn nhất giữa hai đỉnh của đồ thị. Bài toán đó phát biểu dưới dạng
            tổng quát như sau: Cho đồ thị có trọng số $G=(V, E)$, hãy tìm một đường đi ngắn nhất từ đỉnh xuất
            phát $u \in V$ và đỉnh đích $v \in V$. Độ dài của đường đi này ta ký hiệu là $d(u, v)$ và gọi là khoảng
            cách từ $u$ đến $v$. Nếu như không tồn tại đường đi từ $u$ tới $v$ thì ta sẽ đặt khoảng cách đó $=+ \infty$.
            
            Nếu như đồ thị có chu trình âm (chu trình với độ dài âm) thì khoảng cách giữa một số cặp đỉnh nào
            đó có thể không xác định, bởi vì bằng cách đi vòng theo chu trình này một số lần đủ lớn, ta có thể
            chỉ ra đường đi giữa hai đỉnh nào đó trong chu trình này nhỏ hơn bất kỳ một số cho trước nào.
            Lúc này ta xem như không tồn tại đường đi ngắn nhất giữa hai đỉnh.

            Đường đi cơ bản ngắn nhất là đường đi ngắn nhất không có đỉnh được lặp lại quá một lần.
            Nếu như đồ thị không có chu trình âm thì ta có thể chứng minh được rằng một trong những đường
            đi ngắn nhất là đường đi cơ bản.

            \begin{bd}
                Cho một đồ thị có hướng, có trọng số $G=(V, E)$ với trọng số $w: E \rightarrow \mathbb{R}$,
                đặt $p=(v_0, v_1, v_2, \dots, v_k)$ là đường đi ngắn nhất từ đỉnh $v_0$ đến đỉnh $v_k$, bất kỳ $i$ và $j$ với $0 \leq i \leq j \leq k$,
                đặt $p_{ij}=(v_i, v_{i+1}, \dots, v_j)$ là quãng đường con của $p$ từ đỉnh $v_i$ đến đỉnh $v_j$ thì $p_{ij}$ là đường đi ngắn nhất từ đỉnh $v_i$ đến đỉnh $v_j$.

                \label{bd:Shortest-path}
            \end{bd}

            \textbf{Chứng minh:} Ta phân tích $p$ thành các quãng đường nhỏ $v_0 \xrightarrow[]{p_{0i}} v_i \xrightarrow[]{p_{ij}} v_j \xrightarrow[]{p_{jk}} v_k$,
            vì vậy $w(p)=w(p_{0i}) + w(p_{ij}) + w(p_{jk})$. Ta giải sử tồn tại một đường đi $p_{ij}^{\prime}$ từ đỉnh $v_i$ đến đỉnh $v_j$ với trọng số $w(p_{ij}^{\prime}) < w(p_{ij})$.
            Vậy nên quãng đường $v_0 \xrightarrow[]{p_{0i}} v_i \xrightarrow[]{p_{ij}^{\prime}} v_j \xrightarrow[]{p_{jk}} v_k$ với trọng số $w(p_{0i}) + w(p_{ij}^{\prime}) + w(p_{jk}) < w(p)$,
            mâu thuẫn với giả thiết $p$ là quãng đường nhỏ nhất từ $v_0$ đến $v_k$.
            
            Các thuật toán tìm đường đi ngắn nhất thường dựa trên đặc tính là đường đi ngắn nhất giữa hai đỉnh bao gồm các đường đi ngắn nhất giữa các đỉnh trung gian.
            Cấu trúc con tối ưu thường là dấu hiệu để sử dụng các phương pháp như quy hoạch động, tham lam.

            Một số những biến thể của bài toán đường đi ngắn nhất:

            \begin{itemize}
                \item Bài toán đường đi ngắn nhất với đỉnh xuất phát đơn: Tìm một đường đi ngắn nhất từ một đỉnh $s \in V$ cho trước tới tất cả các đỉnh $v \in V, v \neq s$ còn lại.
                \item Bài toán tìm đường đi ngắn nhất giữa một cặp đỉnh: Tìm một đường đi ngắn nhất từ $u$ đến $v$ ($u, v \in V, u \neq v$).
                Nếu ta đang giải bài toán đường đi ngắn nhất với đỉnh xuất phát đơn thì ta cũng đang giải bài toán đường đi ngắn nhất giữa một cặp đỉnh.
                \item Bài toán đường đi ngắn nhất giữa mọi cặp đỉnh: Tìm đường đi ngắn nhất giữa mọi cặp đỉnh $(u, v)$. 
                Ta có thể giải bài toán tìm đường đi ngắn nhất với đỉnh xuất phát đơn và cho từng đỉnh nhưng ta sẽ có một số thuật toán nhanh hơn.
            \end{itemize}
            \item Một số phương pháp (thuật toán) giải bài toán tìm đường đi (ngắn nhất)
            \begin{itemize}
                \item \textbf{Thuật toán Bellman-Ford}:
                
                Thuật toán Bellman-Ford để giải các bài toán đường đi ngắn nhất với đỉnh xuất phát đơn.
                Thuật toán Bellman-Ford được thiết kế dựa trên phương pháp quy hoạch động.

                
                Ta có một nhận xét quan trọng rằng đường đi nhỏ nhất giữa hai đỉnh trong đồ thị (giả sử rằng không tồn tại chu trình âm), thì đường đi này đi qua không quá $\lvert E \rvert - 1$ cạnh.
                Hay nói cách khác không có đỉnh nào được đi qua hai lần trong đường đi ngắn nhất giữa hai đỉnh bất kỳ.
                Để chứng minh, ta giả sử trong một đường đi ngắn nhất giữa hai đỉnh $s, v$, tồn tại một đỉnh được đi qua hai lần hay nói cách khác tồn tại một chu trình với trọng số dương (giả sử rằng không tồn tại chu trình trọng số âm trong đồ thị).
                Ta gọi chu trình này $c = (v_i, v_{i+1}, \dots, v_l)$ với $v_i \equiv v_l$, do $w(c) > 0$, ta hoàn toàn có thể tìm ra một đường đi ngắn hơn bằng cách bỏ chu trình $c$ và chỉ đi qua đỉnh $v_i$.
                Hệ quả là một đường đi ngắn nhất giữa hai đỉnh $s, v$ bất kỳ nếu tồn tại thì không bao giờ đi qua quá $\lvert V \rvert - 1$ đỉnh (không tính đỉnh bắt đầu) hay $\lvert E \rvert - 1$ cạnh.

                Ta gọi $d^{i}(s, v)$ là trọng số quãng đường nhỏ nhất từ đỉnh bắt đầu $s$ đến đỉnh $v$ mà quãng đường này đi không quá $i$ đỉnh tiếp theo từ đỉnh bắt đầu $s$ (không tính điểm bắt đầu $s$).
                Ở bước lặp thứ $i$, ta đã có trọng số độ dài quãng đường nhỏ nhất từ đỉnh $s$ đến đỉnh $v$ đi qua không quá $i - 1$ đỉnh.
                Đường đi từ $s$ đến $v$ không quá $i$ đỉnh được tính bằng cách đi từ $s$ đến một điểm $u$ nào đó qua không quá $i - 1$ đỉnh rồi từ đỉnh $u$ đi trực tiếp đến đỉnh $v$ (nếu $(u, v) \in E$).
                Vì vậy $d^{i}(s, v)$ được chọn là trường hợp tốt nhất trong các trường hợp trên:

                \begin{equation*}
                    d^{i}(s, v) = \min \big( d^{i - 1}(s, v), d^{i - 1}(s, u) + w(u, v) \text{ với } \forall u: (u, v) \in E \big)
                \end{equation*}

                Khi lập trình ta không cần tính $d^{i}(s, v)$ theo $i$ (không cần mảng 2 chiều, 1 chiều để lưu giá trị nhỏ nhất từ đỉnh $s$ đến $v$, một chiều theo $i$),
                ta sẽ tính trực tiếp $d(s, v)$ tại bước tính tối ưu thứ $i$ sẽ được tính toán tối ưu qua các giá trị $d(s, u), \text{ với } \forall u: (u, v) \in E$ tại bước thứ $i - 1$.
                Ta sẽ tiết kiệm được bộ nhớ chỉ phải lưu mảng một chiều thay vì phải lưu mảng hai chiều.

                Để truy vết tìm đường đi ngắn nhất ta tạo một mảng $p$ với $p(v)$ cho biết đỉnh $u$ là đỉnh liền kề trước đỉnh $v$ trong đoạn đường ngắn nhất từ đỉnh xuất phát $s$ đến đỉnh $v$.
                Ban đầu ta khởi tạo $d(s, v) \gets \infty, \thickspace \forall v \in E, v \neq s$, $d(s, s) \gets 0$, $p(v) \gets null, \thickspace v \in V$.
                
                Để phát hiện chu trình âm, ta thực hiện vòng lặp qua các cạnh $(u, v) \in E$ của đồ thị.
                Nếu $d(s, v) > d(s, u) + w(u, v)$, điều này cho thấy đồ thị đã chứa chu trình âm.
                Lý do là vì lập luận ở trên, một đường đi ngắn nhất giữa hai đỉnh bất kỳ sẽ đi qua không quá $\lvert V \rvert - 1$ đỉnh (không tính đỉnh xuất phát $s$),
                ta chỉ cần $\lvert V \rvert - 1$ lần cập nhật để chắc chắn tìm ra đường đi ngắn nhất, nhưng khi ta lặp qua các cạnh thêm lần nữa (lần lặp thứ $\lvert V \rvert$)
                nhưng vẫn tồn tại một đường đi ngắn hơn, điều này chứng tỏ đã có một chu trình âm trong quãng đường từ đỉnh xuất phát $s$ đến đỉnh $v$.

                Ta có sơ đồ mã giả thuật toán Bellman-Ford ở thuật toán \ref{alg:Bellman-Ford}: 

                \begin{algorithm}[h!]
                    \DontPrintSemicolon
                    \KwIn{Đồ thị $G=(V, E)$, ma trận trọng số $w$ với $w(u, v)$ cho biết trọng số giữa hai đỉnh $u, v \in V$, đỉnh xuất phát $s \in V$}
                    \KwOut{Trả về quãng đường ngắn nhất từ $s$ tới các đỉnh $v \in V \backslash \lbrace s \rbrace$ và bảng truy vết $p$}
                    \ForEach{$v \in V$}{
                        $d(s, v) \gets \infty$\;
                        $p(v) \gets null $\;
                    }
                    $d(s, s) \gets 0$\;

                    \For {$i \gets 1$ \textbf{to} $\lvert V \rvert - 1$}{
                        \ForEach{$(u, v) \in E$}{
                            \If {$d(s, v) > d(s, u) + w(u, v)$}{
                                $d(s, v) \gets d(s, u) + w(u, v)$\;
                                $p(v) \gets u$\;
                            }
                        }
                    }

                    \ForEach{$(u, v) \in E$}{
                        \If {$d(s, v) < \infty$ \textbf{and} $d(s, v) > d(s, u) + w(u, v)$}{
                            $d(s, v) \gets -\infty$\;
                            $p(v) \gets u$\;
                        }
                    }
                    \Return{$d, p$}
                    \caption{Thuật toán Bellman-Ford}
                    \label{alg:Bellman-Ford}
                \end{algorithm}

                Để truy vết tìm đường đi ngắn nhất từ đỉnh bắt đầu $s$ đến đỉnh $v$ (giả sử đồ thị $G$ không tồn tại một chu trình âm),
                ta sử dụng mảng $p$ đã được xây dựng ở thuật toán \ref{alg:Bellman-Ford}:

                \begin{algorithm}[h!]
                    \DontPrintSemicolon
                    \KwIn{Mảng truy vết $p$, đỉnh bắt đầu $s$, đỉnh đích $v$}
                    \KwOut{Đường đi từ đỉnh bắt đầu $s$ đến đỉnh đích $v$}
                    
                    $path \gets \lbrack \rbrack$\;
                    $u \gets v$\;
                    \While{$u \neq s$}{
                        $u \gets p(u)$\;
                        $path \gets path.append(u)$\;
                    }
                    $path \gets reverse(path)$\;
                    \Return{$path$}\;
                    \caption{Truy vết thuật toán Bellman-Ford}
                    \label{alg:Bellman-Ford-Trace}
                \end{algorithm}

                Độ phức tạp về chi phí tính toán của thuật toán Bellman-Ford:
                Bước khởi tạo ta thực hiện vòng lặp với $\lvert V \rvert$ bước lặp.
                Bước cập nhật các giá trị tối ưu gồm hai vòng lặp lồng nhau, vòng lặp bên ngoài cần $\lvert V \rvert - 1$ bước lặp.
                Vòng lặp bên trong cần $\lvert E \rvert$ bước lặp. Trong đồ thị đặc, $\lvert E \rvert = \lvert V \rvert^2$, ta có độ phức tạp tính toán của bước cập nhật là $O(\lvert V \rvert^3)$
                Bước kiểm tra có tồn tại chu kỳ âm không ta thực hiện vòng lặp với $\lvert E \rvert$ bước lặp, chi phí của bước kiểm tra này là $O(\lvert V \rvert^2)$.
                Tổng hợp lại ta có độ phức tạp về chi phí tính toán của thuật toán Bellman-Ford là $O(\lvert V \rvert^2)$.
                
                Độ phức tạp về chi phí bộ nhớ: Ta cần lưu hai mảng một chiều $d$ để lưu trọng số quãng đường nhỏ nhất từ đỉnh xuất phát $s$ đến đỉnh $v$ và mảng $p$ với $u = p(v)$ là đỉnh kề trước đỉnh $v$ trong quãng đường ngắn nhất từ đỉnh xuất phát $s$ đến đỉnh đích $v$ với chi phí $O(\lvert V \rvert)$
                \item \textbf{Thuật toán Dijkstra}:

                Thuật toán Dijkstra giải các bài toán đường đi ngắn nhất trên đồ thị có trọng số
                $G=(V, E)$ nhưng yêu cầu tất cả các trọng số của các cạnh phải không âm: $w(u, v) \leq 0 \thickspace \forall \thickspace (u, v) \in E$.
                
                Trong trường hợp trọng số trên các cạnh không âm, thuật toán Dijkstra hiệu quả hơn khá nhiều so với thuật toán Bellman-Ford.
                Đối với thuật toán Bellman-Ford: Với đỉnh $v \in V$, Gọi $d(s, v)$ là độ dài đường đi ngắn nhất từ $s$ đến $v$.
                Thuật toán Bellman-Ford khởi tạo $d(s, s) \gets 0$ và $d(s, v) \gets \infty, \forall v \in V, v \neq s$.
                Ta tối ưu hóa dần các $d(s, v)$ bằng công thưc cập nhật $d(s, v) \gets d(s, u) + w(u, v), \forall (u, v) \in E$.
                Vậy ta dùng đỉnh $u$ để cập nhật $d(s, v)$, nhưng sau nếu ta tối ưu được $d(s, nu)$, ta cũng phải sửa lại nhãn $d(s, v)$ dẫn đến hệ quả $d(s, v)$ phải cập nhật rất nhiều lần.
                Vậy nên ý tưởng của thuật toán Dijkstra là tại mỗi bước ta không xét mọi cặp đỉnh $(u,v)$ để cập nhật $d(s, v)$ qua $d(s, u)$,
                ta chỉ chọn đỉnh $u$ mà khoảng cách $d(s, u)$ không thể tối ưu thêm được nữa.

                Các bước để thực hiện thuật toán Dijkstra:

                \textbf{Bước 1:} Khởi tại. Với đỉnh $v \in V$, gọi $d(s, v)$ là đường đi ngắn nhất từ $s$ đến $v$.
                Ta sẽ tính các $d(s, v)$. Ban đầu $d(s, v)$ được khởi tạo giống như trong thuật toán Bellman-Ford ($d(s, s) \gets 0$, $d(s, v) \gets \infty, \forall v \in V, v \neq s$).
                Ta tạo hai tập hợp $S$ và $Q$. $S$ là tập hợp các đỉnh có mà khoảng cách đã được tối ưu, còn tập $Q$ gồm các đỉnh mà khoảng cách có thể được tối ưu hơn nữa.

                \textbf{Bước 2:} Thực hiện một vòng lặp, mỗi bước lặp gồm 2 thao tác:

                \begin{itemize}
                    \item Trong các đỉnh thuộc tập $Q$, chọn ra đỉnh $u: u = \underset{v \in Q}{\mathrm{argmin} \thickspace d(s, v)}$.
                    Chuyển $u$ sang tập $S$.
                    \item Cập nhật các đỉnh $v \in Q$ theo công thức:
                    \begin{equation*}
                        d(s, v) = \min(d(s, v), d(s, u) + w(u, v))
                    \end{equation*}
                \end{itemize}

                Vòng lặp sẽ kết thức khi tập $Q$ rỗng, hoặc $d(s, v) = \infty \thickspace \forall \thickspace v \in Q$.

                Một điều cần lưu ý rằng vì sao đỉnh $u$ sẽ không thể tối ưu thêm được khoảng cách $d(s, u)$ nữa?
                Ta giả sử $d(s, u)$ còn cố thể tối ưu thêm được nữa vì vậy cần phải tồn tại một đỉnh $t \in Q$ sao cho $d(s, u) > d(s, t) + w(t, u)$.
                Vì $w(t, u) \geq 0$ nên $d(s, u) > d(s, t)$, trái với lựa chọn $u: u = \underset{v \in Q}{\mathrm{argmin} \thickspace d(s, v)}$,
                vậy nên bắt buộc $d(s, u)$ không còn thể tối ưu thêm được nữa.

                Từ các phân tích trên, ta có sơ đồ mã giả của thuật toán Dijkstra được trình bày ở thuật toán \ref{alg:Dijkstra}
                \begin{algorithm}[h!]
                    \DontPrintSemicolon
                    \KwIn{Đồ thị $G=(V, E)$, ma trận trọng số $w$ với $w(u, v)$ cho biết trọng số giữa hai đỉnh $u, v \in V$, đỉnh xuất phát $s \in V$}
                    \KwOut{Trả về quãng đường ngắn nhất từ $s$ tới các đỉnh $v \in V \backslash \lbrace s \rbrace$ và bảng truy vết $p$}

                    \ForEach{$v \in V$}{
                        $d(s, v) \gets \infty$\;
                        $p(v) \gets null $\;
                    }
                    $d(s, s) \gets 0$\;

                    $S \gets \emptyset$\;
                    $Q \gets V$\;

                    \While{$Q \neq \emptyset$ \textbf{or} $\exists v \in Q: d(s, v) < +\infty$}{
                        $u \gets \underset{v \in Q}{\mathrm{argmin}} \thickspace d(s, v)$\;
                        $S \gets S \cup \lbrace u \rbrace$\;
                        $Q \gets Q \backslash \lbrace u \rbrace$\;
                        \ForEach{$v \in Q$ \textbf{and} $v \in adjacent(u)$}{
                            \If{$d(s, v) > d(s, u) + w(u, v)$}{
                                $d(s, v) \gets d(s, u) + w(u, v)$\;
                                $p(v) \gets u$\;
                            }
                        }
                    }
                    \Return{$d, p$}\;
                    \caption{Thuật toán Dijkstra}
                    \label{alg:Dijkstra}
                \end{algorithm}

                Để truy vết đường đi ta sử dụng thuật toán truy vết của thuật toán Bellman-Ford đã được trình bày ở thuật toán \ref{alg:Bellman-Ford-Trace}.
                Độ phức tạp của thuật toán Dijkstra là: Bước khởi tạo cần $\lvert V \rvert$ bước lặp.
                Bước lặp của thuật toán gồm 2 vòng lặp lồng nhau, vòng lặp while bên ngoài trường hợp xấu nhất phải duyệt qua $\lvert V \rvert$ đỉnh.
                Vòng lặp foreach bên trong số lượng sẽ giảm dần sau mỗi vòng lặp while bên ngoài. Nhưng trong trường hợp xấu nhất, đồ thị $G=(V, E)$ là đồ thị kết nối đầy đủ,
                ở bước lặp đầu tiên của vòng lặp while đầu tiên vòng lặp foreach bên trong cần $\lvert V \rvert - 1$ bước lặp.
                Sau mỗi bước lặp while vòng lặp foreach bên trong số lượng bước lặp giảm đi 1.
                Như vậy tổng số vòng lặp bên trong cùng là: $(\lvert V \rvert - 1) + (\lvert V \rvert - 2) + \dots + 1=\dfrac{\lvert V \rvert(\lvert V \rvert - 1)}{2} \sim O(\lvert V \rvert^2)$.
                Độ phức tạp của thuật toán Dijkstra là: $O(\lvert V \rvert + \lvert V \rvert^2)$

                \item \textbf{Thuật toán Floyd-Warshall}:
                
                Thuật toán Floyd-Warshall là thuật toán để giải bài toán tìm đường đi ngắn nhất giữa từng cặp đỉnh.
                Thuật toán Floyd-Warshall được thiết kế dựa trên phương pháp quy hoạch động.
                Ta có thể giải bài toán tìm đường đi ngắn nhất giữa từng cặp đỉnh bằng cách giải bài toán tìm đường đi ngắn nhất với đỉnh xuất phát đơn $\lvert V \rvert$ khả năng chọn đỉnh xuất phát.
                Nhưng với phương pháp này, ta sử dụng thuật toán Bellman-Ford thì độ phức tạp về chi phí tính toán sẽ là $O(\lvert V \rvert^2 \lvert E \rvert)$.
                Nếu với đồ thị đặc $\lvert E \rvert = \lvert V \rvert^2$ thì độ phức tạp về chi phí tính toán sẽ là $O(\lvert V \rvert^4)$.
                Đây là một chi phí rất lớn.

                Thuật toán Floyd-Warshall được thiết kế dựa trên một vài quan sát quan trọng.
                Ta đánh số các đỉnh của đồ thị $G$ là $V=\lbrace 1, 2, \dots, n \rbrace$, ta chọn ra một tập con các đỉnh
                $\lbrace 1, \dots, k$ sao cho $1 \leq k \leq n$. Bất kỳ một cặp đỉnh $i, j \in V$,
                ta xét tất cả các đường đi từ $i$ đến $j$ mà các đỉnh trung gian có thể được chọn trong tập $\lbrace 1, \dots, k \rbrace$,
                ta đặt $p$ là đường đi có trọng số nhất trong các đường đi này.
                Thuật toán Floyd-Warshall khai thác mối quan hệ giữa $p$ và các đường đi nhỏ nhất từ $i$ đến $j$ với các đỉnh trung gian được chọn từ tập $\lbrace 1, \dots, k - 1 \rbrace$.
                Mối quan hệ này phụ thuộc vào đỉnh $k$ có phải là một đỉnh trung gian của $p$ hay không.

                \begin{itemize}
                    \item \textbf{Trường hợp 1:} Nếu đỉnh $k$ không là đỉnh trung gian của đường đi $p$ thì tất cả các đỉnh trung gian của $p$ thuộc về tập $\lbrace 1, 2, \dots, k-1 \rbrace$.
                    Vì vậy mà đường đi ngắn nhất từ đỉnh $i$ đến đỉnh $j$ với tập đỉnh trung gian $\lbrace 1, 2, \dots, k - 1 \rbrace$ cũng là đường đi ngắn nhất từ đỉnh $i$ đến đỉnh $j$ với tập đỉnh trung gian là $\lbrace 1, 2, \dots, k \rbrace$.
                    
                    \item \textbf{Trường hợp 2:} Nếu $k$ là đỉnh trung gian của đường đi $p$, ta tách $p$ thành hai đường $i \xrightarrow[]{p_1} k \xrightarrow[]{p_2} j$.
                    Theo bổ đề \ref{bd:Shortest-path}, $p_1$ là đường đi ngắn nhất từ $i$ đến $k$ với tập đỉnh trung gian $\lbrace 1, 2, \dots, k - 1 \rbrace$.
                    Tương tự, $p_2$ cũng là đường đi ngắn nhất từ đỉnh $k$ đến đỉnh $j$ với tập đỉnh trung gian $\lbrace 1, 2, \dots, k - 1 \rbrace$.
                \end{itemize}
                
                Từ các quan sát trên, ta xây dựng được công thức truy hồi về đường đi nhỏ nhất từ đỉnh $i$ đến đỉnh $j$ với tập đỉnh trung gian $\lbrace 1, 2, \dots, k \rbrace$.
                Ta đặt trọng số đường đi nhỏ nhất từ đỉnh $i$ đến đỉnh $j$ với tập đỉnh trung gian $\lbrace 1, 2, \dots, k \rbrace$ là $d^{k}(i, j)$. 
                Với $k = 0$, một đường đi từ đỉnh $i$ đến đỉnh $j$ với không có đỉnh trung gian nào thì đường đi ngắn nhất chính là trọng số của cạnh nối trực tiếp giữa đỉnh $i$ và đỉnh $j$ (nếu có).
                Ta có công thức truy hồi tính $d^k(i, j)$:

                \begin{equation*}
                    d^{k}(i, j) = \begin{cases} w(i, j) \text{ nếu } k = 0 \\ \min \big \lbrace d^{k-1}(i, j), d^{k-1}(i, k) + d^{k-1}(k, j)  \big\rbrace \text{ nếu } k \geq 1 \end{cases}
                \end{equation*}

                Để truy vết đường đi ngắn nhất giữa các cặp đỉnh, ta sử dụng ma trận $P$ với $P^{k}(i, j)$ có ý nghĩa là trong quãng đường ngắn nhất từ đỉnh $i$ đến đỉnh $j$ với tập đỉnh trung gian là $\lbrace 1, 2, \dots, k \rbrace$ thì điểm liền trước đỉnh $j$ là đỉnh nào.
                Ma trận $P$ được xây dựng trong quá trình thực hiện thuật toán Floyd-Warshall, ban đầu ma trận $P$ được khởi tạo:

                \begin{equation*}
                    P^{0}(i, j) = \begin{cases} null \text{ nếu } i = j \text{ hoặc } w(i, j) = \infty \\ i \text{ nếu } i \neq j \text{ hoặc } w(i, j) < \infty \end{cases}
                \end{equation*}

                Với $k \geq 1$, nếu $k$ là một đỉnh trung gian trong đường đi ngắn nhất từ đỉnh $i$ đến đỉnh $j$ $i \xrightarrow[]{p_1} k \xrightarrow[]{p_2} j (k \neq j)$,
                ta chọn điểm liền trước của $j$ trên đường đi ngắn nhất từ đỉnh $i$ đến đỉnh $j$ cùng với đỉnh là đỉnh liền trước đỉnh $j$ trong quãng đường ngắn nhất từ đỉnh $k$ đến đỉnh $j$ với tập đỉnh trung gian là $\lbrace 1, 2, \dots, k - 1 \rbrace$.
                Ta có công thức tổng quát tính các phần tử của ma trận $P (k \geq 1)$:

                \begin{equation*}
                    P^{k}(i, j) = \begin{cases} P^{k-1}(k, j) \text{ nếu } d^{k-1}(i, j) > d^{k-1}(i, k) + d^{k-1}(k, j) \\
                                                P^{k-1}(i, j) \text{ nếu } d^{k-1}(i, j) \leq d^{k-1}(i, k) + d^{k-1}(k, j) \end{cases}
                \end{equation*}

                Ta cũng có thể tìm một chu trình âm trong đồ thị bằng cách kiểm tra các phần tử trên đường chéo chính của ma trận $d^{n}$.
                Nếu có ít nhất một phần tử âm, ta kết luận rằng đồ thị tồn tại chu trình âm.
                
                Một điều cần lưu ý là trong quá trình l ập trình, ta không còn cần tính $d^{k}(i, j)$ theo $k$ nữa mà sẽ tính trực tiếp $d(i, j)$.
                $d(i, j)$ tại bước tính tối ưu thứ $k$ sẽ được tính toán để tối ưu qua các giá trị $d(i, j), d(i, k), d(k, j)$ tại bước thứ $k-1$.
                Như vậy ta sẽ tiết kiệm được bộ nhớ chỉ phải lưu mảng hai chiều thay vì phải lưu mảng 3 chiều.

                Ta có mã giả thuật toán Floyd-Warshall được trình bày ở thuật toán \ref{alg:Floyd-Warshall}:

                \begin{algorithm}[h!]
                    \DontPrintSemicolon
                    \KwIn{Đồ thị $G=(V, E)$ được biểu diễn bằng ma trận kề $w$}
                    \KwOut{Khoảng cách nhỏ nhất giữa các cặp đỉnh được biểu diễn bằng ma trận $d$ và ma trận truy vết $P$}

                    \For {$i \gets 1$ \textbf{ to } $n$} {
                        \For {$j \gets 1$ \textbf{ to } $n$} {
                            $d(i, j) \gets w(i, j)$\;
                            \If{$i = j$ \textbf{ or } $w(i, j) = \infty$} {
                                $P(i, j) \gets null$\;
                            } \Else{
                                $P(i, j) \gets i$\;
                            }
                        }
                    }

                    \For {$k \gets 1$ \textbf{ to } $n$} {
                        \For {$i \gets 1$ \textbf{ to } $n$} {
                            \For {$k \gets 1$ \textbf{ to } $n$} {
                                \If {$d(i, j) > d(i, k) + d(k, j)$} {
                                    $d(i, j) \gets d(i, k) + d(k, j)$\;
                                    $P(i, j) \gets P(k, j)$\;
                                }
                            }
                        }
                    }

                    \For {$i \gets 1$ \textbf{ to } $n$} {
                        \If{$d(i, i) < 0$} {
                            print("Phát hiện chu trình âm")\;
                            break\;
                        }
                    }
                    \Return{$d, P$}\;
                    \caption{Thuật toán Floyd-Warshall}
                    \label{alg:Floyd-Warshall}
                \end{algorithm}

                Thuật toán Floyd-Warshall gồm 3 công việc riêng biệt:
                \begin{itemize}
                    \item Quá trình khởi tạo gồm 2 vòng lặp for lồng nhau có độ phức tạp về chi phí tính toán là $O(n^2)$.
                    \item Thực hiện cập nhật các đường đi tối ưu và lưu vết gồm 3 vòng lặp lồng nhau có độ phức tạp về chi phí tính toán là $O(n^2)$.
                    \item Quá trình kiểm tra chu trình âm, trường hợp xấu nhất cần $n$ bước kiểm tra, độ phức tạp về chi phí tính toán là $O(n)$.
                \end{itemize}
                Như vậy độ phức tạp về chi phí tính toán của thuật toán Floyd-Warshall là $O(n^3)$.
                Độ phức tạp về chi phí bộ nhớ là cần lưu hai mảng hai chiều $d$ và $P$ là $O(n^2)$.

                Để đưa ra được đường đi ngắn nhất giữa hai đỉnh $s$ và $v$, ta sử dụng thuật toán tương đối giống với thuật toán truy vết của thuật toán Bellman-Ford.
                Mã giả của thuật toán đưa ra đường đi ngắn nhất được trình bày ở thuật toán \ref{alg:Floyd-Warshall-Trace}.

                \begin{algorithm}[h!]
                    \KwIn{Mảng truy vết $P$, đỉnh bắt đầu $s$, đỉnh đích $j$}
                    \KwOut{Đường đi ngắn nhất từ đỉnh bắt đầu $s$ đến đỉnh đích $v$}

                    $path \gets \lbrack \rbrack$\;
                    
                    \While{$v \neq s$}{
                        $v \gets P(s, v)$\;
                        $path \gets path.append(v)$\;
                    }
                    $path \gets reverse(path)$\;
                    \Return{$path$}\;
                    \caption{Truy vết thuật toán Floyd-Warshall}
                    \label{alg:Floyd-Warshall-Trace}
                \end{algorithm}
            \end{itemize}
        \end{itemize}
    \end{loigiai}
    \newpage

    \printbibliography[title={TÀI LIỆU THAM KHẢO}]

\end{document}